\unit{Exposé 221 -- Hilbert schemes}
\sid{221}
\src{249}
\seminar{13th year, 1960/61, no. 221\\May 1961}
\paperhead{CONSTRUCTION TECHNIQUES AND EXISTENCE THEOREMS IN ALGEBRAIC GEOMETRY\\
IV. HILBERT SCHEMES}%
  {by Alexander GROTHENDIECK}

\subsection*{Introduction}
\addcontentsline{toc}{subsection}{Introduction}

The techniques presented in [2], I and II, were for the most part independent
of any projectivity hypothesis on the schemes under consideration.
Unfortunately, they do not at present make it possible to solve the existence
problems posed in Exposé II of [2]. In the present exposé and the next, we shall
solve these problems under projectivity hypotheses. The techniques employed
are typically projective and make practically no use of the results of I and
II of [2]. Here we shall construct the “Hilbert schemes” intended to replace
the use of Chow coordinates, as stated in [2], II, no. 2. In the next exposé,
the theory of passage to the quotient in schemes developed in III, together
with the theory of Hilbert schemes, will enable us, for example, to construct
Picard schemes (defined in [2], II, no. 3) under fairly general conditions.

In summary, one may say that a fairly satisfactory technique of projective
constructions is now available, except that there is still
lacking\footnote{See the addendum at the end of the exposé.} a theorem on
passage to the quotient by groups such as the projective group, acting
“without fixed points” (cf. [2], III, no. 8). The situation even appears
somewhat better in analytic geometry (when one restricts to studying analytic
spaces “projective” over a given analytic space), since for analytic spaces
the difficulty of passing to the quotient by a group acting nicely disappears.
Moreover, in algebraic geometry as in analytic geometry, it remains to develop
a construction technique valid without projectivity hypotheses.

\fsection{1}{Bounded families of sheaves: permanence properties}

Let $k$ be a field and $X$ a $k$-prescheme, of finite type for simplicity. For
every extension $K/k$ one thus obtains a $K$-prescheme
$X_K=X\otimes_k K$. If $\mathcal F$ is a coherent sheaf on $X_K$ and $K'$ is
an extension of $K$, then $\mathcal F\otimes_K K'=\mathcal F_{K'}$ is a
quasi-coherent sheaf on $X_K\otimes_K K'=X_{K'}$. Now, if $K$ and $K'$ are
any two extensions of $k$, $\mathcal F$ is a quasi-coherent sheaf on $X_K$,
and $\mathcal F'$
\src{250}
is a quasi-coherent sheaf on $X_{K'}$, we say that $\mathcal F$ and
$\mathcal F'$ are \emph{equivalent} if there exist $k$-homomorphisms from
$K$ and $K'$ into a common extension $K''$ of $k$ such that
$\mathcal F_{K''}$ and $\mathcal F'_{K''}$ are isomorphic on $X_{K''}$. This
is an equivalence relation, and we shall be interested in equivalence classes
of sheaves under this relation and in sets of equivalence classes. Note that
if $X$ is of finite type over $k$,\footnote{The printed text has $X_0$ here;
only $X$ is defined in this passage.} every class of coherent sheaves can be
defined by a coherent sheaf on some $X_K$, where $K$ is an extension of finite
type over $k$. In defining classes of coherent sheaves one may therefore
restrict to algebraically closed extensions of $k$; one may also restrict to a
fixed algebraically closed extension $\Omega$ of $k$ of infinite
transcendence degree. Two coherent sheaves $\mathcal F$ and $\mathcal F'$ on
$X_\Omega$ are then equivalent if and only if there exists a
$k$-automorphism $\sigma$ of $\Omega$ such that
$\mathcal F\otimes_\Omega(\Omega,\sigma)$ is isomorphic to $\mathcal F'$.
There is thus a bijective correspondence between the classes of coherent
sheaves under either definition.

Let $E$ and $E'$ be two sets of classes of coherent sheaves on $X$. Consider
the classes of all sheaves of the form $\mathcal F\otimes\mathcal F'$, where
$\mathcal F$ and $\mathcal F'$ are coherent sheaves on the same $X_K$, the
class of $\mathcal F$ belonging to $E$ and that of $\mathcal F'$ to $E'$. One
thereby obtains a set of classes of coherent sheaves, denoted $E\otimes E'$.
The sets $\operatorname{Tor}_i(E,E')$, etc., are defined in the same way. More
generally, to every function $U$ assigning to each sequence
$\mathcal F_1,\ldots,\mathcal F_n$ of $n$ coherent sheaves on a common $X_K$
a set $U(\mathcal F_1,\ldots,\mathcal F_n)$ of coherent sheaves on $X_K$, and
having the evident compatibility with sheaf isomorphisms and inverse images
under base change, we associate a function, denoted by the same symbol $U$,
which assigns to every sequence of $n$ sets $E_1,\ldots,E_n$ of classes of
coherent sheaves a set $U(E_1,\ldots,E_n)$ of classes of coherent sheaves.

Our purpose in this section is to define certain sets of classes of sheaves,
called \emph{bounded}, and to show that the most common operations $U$, when
applied to bounded sets, again give bounded sets.

Let $X$ be a prescheme of finite type over a noetherian $S$. For every
$s\in S$, the fiber $X_s$ is a prescheme of finite type over $k(s)$, and we
shall consider classes of coherent sheaves on $X_s$ in the preceding sense.
This gives meaning to the phrase “class of coherent sheaves on a fiber of
$X/S$,” and to analogous phrases. Likewise, proceeding separately on each
fiber, one may still consider operations such as $E\otimes E'$, etc., which
assign to systems of sets of
\src{251}
classes of coherent sheaves on the fibers of $X/S$ another set of classes of
coherent sheaves on the fibers of $X/S$.

\result{DEFINITION}{1.1} Let $E$ be a set of classes of coherent sheaves on the
fibers of $X/S$. We say that $E$ is \emph{bounded} if there exist a prescheme
$S'$ of finite type over $S$ and a coherent sheaf $\mathcal F'$ on
$X'=X\times_S S'$ such that $E$ is contained in the set of classes of sheaves
on the fibers of $X/S$ defined by $\mathcal F'$.

By definition, the latter assigns to an $s\in S$ the classes of the sheaves
$\mathcal F'\otimes_{S'}k(s')$, where $s'$ ranges over the points of $S'$ above
$s$ (so that $k(s')$ is an extension of $k(s)$ and
$X_s\otimes_{k(s)}k(s')$ is identified with the fiber
$X'\otimes_{S'}k(s')=X'_{s'}$ of $X'$ at $s'$). Thus bounded families may be
described as those contained in an \emph{algebraic family} of coherent sheaves
parametrized by an $S'$ of finite type over $S$.

A finite union of bounded families is bounded (take the disjoint-union
prescheme of the parameter preschemes $S_i$ defining the dominating algebraic
families). A base change $T\longrightarrow S$ transforms a bounded family
relative to $X/S$ into a bounded family relative to $X_T/T$, and the converse
holds if $T\longrightarrow S$ is surjective (or, more generally, if its image
contains the $s$ that actually occur in the given family $E$ for $X/S$). This
reduces, in principle, the determination of bounded families to the case in
which $S$ is the spectrum of an algebra of finite type over the ring of
integers $\mathbf Z$.

If $E$ and $E'$ are bounded families of classes of sheaves relative to $X/S$,
then $E\otimes E'$ is also bounded. Indeed, suppose $E$ and $E'$ are dominated,
respectively, by the algebraic families defined by $T\longrightarrow S$ and
$\mathcal F$ on $X_T$, and by $T'\longrightarrow S$ and $\mathcal F'$ on
$X_{T'}$. Then $E\otimes E'$ is dominated by the algebraic family defined by
$T''=T\times_S T'\longrightarrow S$ and the sheaf $\mathcal F''$ on
$X_{T''}$ obtained as the tensor product of the inverse images of $\mathcal F$
and $\mathcal F'$ on $X_{T''}$. This argument is valid only because the functor
$\mathcal F\otimes\mathcal F'$ in $\mathcal F,\mathcal F'$ is right exact and
therefore commutes with base extension (and in particular with passage to the
fibers). It is not directly applicable to local operations such as
$\operatorname{Tor}_i(E,E')$, $\operatorname{Hom}(E,E')$, and
$\operatorname{Ext}^i(E,E')$. One can nevertheless show that these operations
again transform bounded sets\footnote{The printed French omits the determiner \emph{des}; the English wording was already idiomatic.} into bounded sets, proceeding as for
$E\otimes E'$ but also using results of the following kind (all contained in
[3], IV, 6.11): a bounded family $E$ is always dominated by an algebraic
family defined by a coherent sheaf $\mathcal F$ on an $X_T$ ($T$ of finite
type over $S$) that is flat relative to $T$. (One “cuts into pieces”\footnote{The printed French has the malformed phrase « coupe au morceaux »; the English wording already expresses the corrected « coupe en morceaux ».} the
initial parameter space.) Such
\src{252}
flatness properties for suitable sheaves indeed ensure that operations such as
$\operatorname{Tor}_i(\mathcal F,\mathcal F')$ commute with arbitrary base
change. The same method applies to operations of a global nature: direct
images and derived direct images of coherent sheaves under proper morphisms,
global $\operatorname{Ext}$ relative to proper morphisms (cf. [5], III,
\S~6), etc. All these operations transform bounded families of sheaves into
bounded families of sheaves. (N.B. -- Here the preschemes on which the various
sheaves are taken may change under operations of the kind considered.)

The following two statements are proved essentially by the same flatness
technique; for primary decomposition on the fibers of a morphism of finite
type, see in particular [5], IV.

\result{PROPOSITION}{1.2} Let $E$ and $E'$ be bounded sets of classes of
sheaves on the fibers of $X/S$, where $X$ is assumed proper over $S$. Then:

\medskip
\noindent (i) The family of kernels, cokernels, and images of homomorphisms
$\mathcal F\longrightarrow\mathcal F'$, where the class of $\mathcal F$
belongs to $E$ and that of $\mathcal F'$ belongs to $E'$, is bounded.

\smallskip
\noindent (ii) The family of sheaves $\mathcal F''$ that are extensions of a
$\mathcal F$ by a $\mathcal F'$, where the class of $\mathcal F$ belongs to
$E$ and that of $\mathcal F'$ belongs to $E'$, is bounded.

After a suitable base change, one may suppose that $E$ and $E'$ are defined,
respectively, by coherent sheaves $\mathcal G$ and $\mathcal G'$ on an
$X_T/T$, where $T$ is of finite type over $S$. One may further suppose that
certain flatness hypotheses hold, implying that the formation of the sheaves
$f_{T*}(\operatorname{Hom}_{X_T}(\mathcal G,\mathcal G'))$ and
$\operatorname{Ext}^{1}_{f_T}(\mathcal G,\mathcal G')$ commutes with base
change for an arbitrary morphism $T'\longrightarrow T$. One may moreover
suppose that these coherent sheaves on $T$ are locally free. Let $T_0$ and
$T_1$ be the vector bundles over $T$ whose sheaves of germs of sections are,
respectively, the preceding sheaves. One then canonically defines a
homomorphism $\mathcal G_{T_0}\longrightarrow\mathcal G'_{T_0}$ of coherent
sheaves on $X_{T_0}$, and an extension\footnote{The printed text reverses the two end terms of the universal extension; the class under consideration lies in $\operatorname{Ext}^1(\mathcal G,\mathcal G')$.}
\[
0\longrightarrow\mathcal G'_{T_1}\longrightarrow\mathcal G''
\longrightarrow\mathcal G_{T_1}\longrightarrow0
\]
of coherent sheaves on $X_{T_1}$ having an evident universal property. The
second sheaf defines an algebraic family dominating the family considered in
(ii). The same is true of the kernel, cokernel, and image of the homomorphism
\src{253}
above, whose consideration proves (i) (provided that the cokernel is assumed
flat relative to $T_0$, a case to which one again reduces by cutting $T_0$ into
pieces\ldots).

\result{PROPOSITION}{1.3} Let $E$ be a bounded family of classes\footnote{The printed French has the singular « classe »; the English was already plural.} of sheaves on
the fibers of $X/S$. Then the classes of the structure sheaves of
$(\operatorname{supp}\mathcal F)_{\mathrm{red}}$, where $\mathcal F$ is a
coherent sheaf on an $X_K$, with $K$ algebraically closed, whose class belongs
to $E$, form a bounded family.

Here $(\operatorname{supp}\mathcal F)_{\mathrm{red}}$ denotes the support of
$\mathcal F$ endowed with the induced reduced structure; that is, its
structure sheaf is the quotient of $\mathcal O_{X_K}$ by the largest sheaf of
ideals defining $\operatorname{supp}\mathcal F$. One can prove the analogue of
(1.3) for the sheaves canonically deduced from $\mathcal F$ by the theory of
primary decomposition, for example the $\mathcal F/\mathcal F_i$, where the
$\mathcal F_i$ are the subsheaves of $\mathcal F$ primary for the components
of the support of $\mathcal F$ and minimal with this property; or the
$\mathcal O_{X_K}/\mathfrak p$, where $\mathfrak p$ is a sheaf of prime ideals
associated with $\mathcal F$; or the $\mathcal O_{X_K}/\mathfrak q$, where
$\mathfrak q$ is a sheaf of primary ideals associated with a component of the
support of $\mathcal F$ (the reference field being algebraically closed).

\fsection{2}{Bounded families and the Hilbert polynomial}

In what follows, we suppose that $X$ is projective over $S$ and is equipped
with a very ample sheaf denoted $\mathcal O_X(1)$. For every extension $K$ of
a residue-field extension $k(s)$ of a point $s$ of $S$, we consider on $X_K$
the corresponding sheaf $\mathcal O_{X_K}(1)$, which is again very ample.

To every coherent sheaf $\mathcal F$ on $X_K$ we associate the function
\[
P_{\mathcal F}(n)=\text{the Euler--Poincaré characteristic of $\mathcal F(n)$
on $X_K$},
\]
which is a polynomial in the integer $n$, called the \emph{Hilbert polynomial}
of $\mathcal F$. For large values of $n$, $P(n)$ is simply the dimension of
$H^0(X_K,\mathcal F(n))$ over $K$, since the $H^i(X_K,\mathcal F(n))$ vanish
for $i>0$ and $n$ large.

Now suppose that $\mathcal F$ is a coherent sheaf on $X$, flat over $S$. Then
the Hilbert polynomials of the sheaves $\mathcal F_s$ induced on fibers $X_s$
belonging to the same connected component of $S$ are all equal [5], III,
\S~7. It follows (without a flatness hypothesis) that the set of Hilbert
polynomials of the sheaves $\mathcal F_s$, $s\in S$, is finite for every
coherent sheaf $\mathcal F$ on $X$.
\src{254}
Recall also that if $\mathcal F$ is a coherent sheaf on $X$, it is isomorphic
to a quotient sheaf of a sheaf $\mathcal O_X(-n)^N$, for $n$ and $N$
sufficiently large. Thus the sheaves $\mathcal F_s$ induced on the fibers are
also quotients of the sheaf $\mathcal O(-n)^N$ on the fiber.

These two observations give the “only if” part of the following theorem:

\result{THEOREM}{2.1} Let $X$ be projective over a noetherian $S$, with
$\mathcal O_X(1)$ very ample on $X$ relative to $S$. Let $E$ be a set of
classes of sheaves on the fibers of $X/S$. For $E$ to be bounded, it is
necessary and sufficient\footnote{The printed exposé says only “necessary”;
the 1962 erratum to Exposé 221 prescribes “necessary and sufficient.”} that it
satisfy the following conditions:

\medskip
\noindent (a) There exists a coherent sheaf $\mathcal L$ on $X$ (which may be
assumed to have the form $\mathcal O_X(-n)^N$) such that $E$ is contained in
the family of classes of coherent quotient sheaves of sheaves of the form
$\mathcal L_K$.

\smallskip
\noindent (b) The Hilbert polynomials $P_{\mathcal F}$ of the sheaves
$\mathcal F$ whose class belongs to $E$ are members of one fixed finite set of
polynomials.

It remains to prove the “if” part, which will be a special case of a more
precise result. For every coherent Module $\mathcal F$ on a prescheme of finite
type over a field $K$, and every integer $r$, let $N_r$ be the sub-Module of
$\mathcal F$ whose sections on an open set are those sections of $\mathcal F$
on that open set whose support has dimension $<r$. Thus $N_r=\mathcal F$ for
$r>\dim\operatorname{supp}\mathcal F$, while $N_r=0$ for $r\leq0$, and one
obtains a finite increasing filtration of $\mathcal F$ whose factors
$N_{r+1}/N_r$\footnote{The printed text reverses the quotient as
$N_r/N_{r+1}$, which is incompatible with the increasing filtration just
defined.} have as their associated prime cycles exactly the prime cycles
associated with $\mathcal F$ that have dimension $r$. Put
\[
\mathcal F_{(r)}=\mathcal F/N_r.
\]
Then the associated prime cycles of $\mathcal F_{(r)}$ are exactly those prime
cycles associated with $\mathcal F$ that have dimension $\geq r$; in
particular, it equals $\mathcal F$ if and only if the prime cycles associated
with $\mathcal F$ have dimension $\geq r$. With this notation:

\result{THEOREM}{2.2} Under the hypotheses of (2.1), let $s$ be an integer and
suppose that $E$ satisfies condition (a) and the following weakened form of
(b):

\medskip
\noindent $(b_s)$ The Poincaré polynomials $P_{\mathcal F}$ of the sheaves
$\mathcal F$ whose class belongs to $E$ have bounded coefficients in degrees
$\geq s-1$.

Under these conditions, the sheaves $\mathcal F_{(s)}$ (as the class of
$\mathcal F$ ranges through $E$) form a bounded family. Moreover, the
coefficients of degree $s-2$ of the $P_{\mathcal F}$ are bounded below.
\src{255}
Thus:

\result{COROLLARY}{2.3} Suppose that the sheaves $\mathcal F$ whose class
belongs to $E$ are such that all their associated prime cycles have dimension
$d$ satisfying $s\leq d\leq r$. Then in (2.1) (b), it is enough to consider
the coefficients of the $P_{\mathcal F}$ in degrees from $s-1$ through $r$.

The remainder of this section is devoted to an outline of the proof of (2.2).
The key lemmas are the following two, the first of which is well known (and
summarizes the useful mathematical content of Chow coordinates):

\result{LEMMA 2.4 (CHOW)}{} Consider the structure sheaves of subschemes $Y$ of
fibers $X_K$ ($K$ ranging over algebraically closed extensions of residue
fields of $S$), where $Y$ is reduced and all its components have the same
dimension $r$ (with $\mathcal O_Y$ regarded as a quotient sheaf of
$\mathcal O_X$). If the degrees of the $Y$ are bounded, the $Y$ form a bounded
family.

Here the degree $a$ of $Y$ can most conveniently be defined by the coefficient
of the leading term of
\[
P_{\mathcal O_Y}=a\frac{n^r}{r!}+\cdots.
\]

\result{LEMMA}{2.5} Let $\mathcal L$ be a coherent sheaf on $X$, and let $E$
be a set of classes of quotient sheaves $\mathcal F$ of sheaves $\mathcal L_K$
($K$ a residue-field extension of $S$). Suppose that the fibers of $X$ over
$S$ have dimension $\leq r$, and write
\[
P_{\mathcal F}(n)=a_{\mathcal F}\frac{n^r}{r!}
+b_{\mathcal F}\frac{n^{r-1}}{(r-1)!}
+\text{terms of degree $<r-1$}.
\]
Then the coefficient $a_{\mathcal F}$ is bounded, while $b_{\mathcal F}$ is
bounded below. If $b_{\mathcal F}$ is bounded, then the family of the
$\mathcal F_{(r)}$ is bounded.

\emph{Proof.} -- After replacing $S$ by a union of subschemes of $S$ covering
$S$, one may suppose there is a finite morphism
$f:X\longrightarrow\mathbf P^r_S$ such that $\mathcal O_X(1)$ is isomorphic to
the inverse image of $\mathcal O_{\mathbf P^r_S}(1)$. Hence for every coherent
sheaf $\mathcal F$ on $X$ one has
$P_{\mathcal F}=P_{f_*(\mathcal F)}$. On the other hand, one checks readily
(by the technique of the preceding section) that a set of sheaves $\mathcal F$
on $X$ is bounded if and only if the set of $f_*(\mathcal F)$ is bounded.
Finally,
\[
f_*(\mathcal F)_{(r)}=f_*(\mathcal F_{(r)}).
\]
\src{256}
This reduces us to the case $X=\mathbf P^r_S$. Moreover, one may suppose that
$\mathcal L=\mathcal O_{\mathbf P^r_S}(k)^s$ for suitable $k$ and $s$. The
coefficient $a_{\mathcal F}$ satisfies
\[
0\leq a_{\mathcal F}\leq s
\]
and is therefore bounded. It is equivalent to say that the coefficient of
$n^{r-1}$ in the $P_{\mathcal F}(n)$ is bounded below (respectively, bounded),
or to say the same for
$P_{\mathcal F}(n-k)=P_{\mathcal F(-k)}(n)$. We are thus reduced to the case
\[
\mathcal L=\mathcal O_{\mathbf P^r_S}^{s}.
\]

Consider the exact sequence
\[
0\longrightarrow N_r\longrightarrow\mathcal F
\longrightarrow\mathcal F_{(r)}\longrightarrow0,
\]
whence
\[
P_{\mathcal F}=P_{\mathcal F_{(r)}}+P_{N_r}.
\]
Since the coefficient of $n^{r-1}$ in $P_{N_r}$ is positive (because
$\dim\operatorname{supp}N_r\leq r-1$), we have
\[
b_{\mathcal F_{(r)}}\leq b_{\mathcal F}.
\]
Thus, in proving the assertions of the lemma, we may replace $\mathcal F$ by
$\mathcal F_{(r)}$; that is, we may assume that the quotient sheaves
$\mathcal F$ of $\mathcal L$ under consideration are torsion-free.

Since $\mathbf P^r_K$ is normal, it follows that $\mathcal F$ is locally free
of rank $a=a_{\mathcal F}$ on an open set $U=\mathbf P^r_K-Y$, where $Y$ has
codimension $\geq2$. Thus $\bigwedge^a\mathcal F$ is a sheaf on
$\mathbf P^r_K$ whose restriction to $U$ is invertible; hence (since
$\mathbf P^r_K$ is regular and $Y$ has codimension $\geq2$) it is isomorphic
to the restriction of an invertible sheaf on $\mathbf P^r_K$, determined up
to isomorphism. The latter has the form $\mathcal O_{\mathbf P^r_K}(d)$ for a
well-determined integer $d$. Since $\bigwedge^a\mathcal F$ is a quotient of
$\bigwedge^a\mathcal O_{\mathbf P^r_K}^{s}\simeq
\mathcal O_{\mathbf P^r_K}^{N}$\footnote{The printed text replaces the
exponent $s$ here by an undefined $n$; $s$ is the rank of the free sheaf
surjecting onto $\mathcal F$.}, where $N=\binom{s}{a}$, it has $N$ canonical
sections, and these therefore define sections of
\src{257}
$\mathcal O_{\mathbf P^r_K}(d)$ on $U$ that are restrictions of sections $s_i$
($1\leq i\leq N$) of $\mathcal O_{\mathbf P^r_K}(d)$ (since
$\mathbf P^r_K$ is normal and $Y$ has codimension $\geq2$). These sections
generate $\mathcal O_{\mathbf P^r_K}(d)$ at the points of $U$, so they are not
all zero; this implies $d\geq0$. On the other hand, an easy calculation gives
\[
b_{\mathcal F}=a_{\mathcal F}(r+1)/2+d.
\]
In particular, this shows that $b_{\mathcal F}>0$, so $b_{\mathcal F}$ is
bounded below. It is bounded if and only if $d$ is bounded; let us show that
in this case $\mathcal F$ belongs to a bounded family. We may suppose that
$a_{\mathcal F}$ and $b_{\mathcal F}$ are fixed, say equal to $a$ and $b$, so
that $d$ is fixed. Knowledge of the $N$ sections $s_i$ of
$\mathcal O_{\mathbf P^r_K}(d)$, that is, of a homomorphism
\[
s:\bigwedge^a\mathcal L_K\longrightarrow\mathcal O_{\mathbf P^r_K}(d),
\]
allows one to recover $\mathcal F$ as the coimage of the corresponding
composite homomorphism
\[
\mathcal L_K\longrightarrow
\operatorname{Hom}\!\left(\bigwedge^{a-1}\mathcal L_K,
\bigwedge^a\mathcal L_K\right)\longrightarrow
\operatorname{Hom}\!\left(\bigwedge^{a-1}\mathcal L_K,
\mathcal O_{\mathbf P^r_K}(d)\right),
\]
where the first is the canonical homomorphism arising from the exterior
product and the second is deduced from $s$. The conclusion now follows from
(1.2) (i).

Combining the preceding two lemmas proves:

\result{LEMMA}{2.6} Under the preliminary hypotheses of (2.1), suppose that for
every $\mathcal F$ one has
\[
P_{\mathcal F}(n)=a_{\mathcal F}\frac{n^r}{r!}
+b_{\mathcal F}\frac{n^{r-1}}{(r-1)!}
+\text{terms of degree $<r-1$},
\]
and that the coefficients $a_{\mathcal F}$ are bounded. Then the coefficients
$b_{\mathcal F}$ are bounded below. If the $b_{\mathcal F}$ are bounded, then
the $\mathcal F_{(r)}$ form a bounded family.

We may suppose that the base fields $K$ for the sheaves $\mathcal F$ are
algebraically closed. Endow each
$\operatorname{supp}\mathcal F_{(r)}$, the union of the components of
$\operatorname{supp}\mathcal F$ of dimension $r$, with the induced reduced
structure. Then the degrees of the $\operatorname{supp}\mathcal F_{(r)}$ are
bounded above by $a$, so by (2.4) the $\operatorname{supp}\mathcal F_{(r)}$
form a bounded set. Moreover, for every component of
$\operatorname{supp}\mathcal F_{(r)}$, the length of $\mathcal F_{(r)}$ along
that component is $\leq a$. Hence if $\mathcal J_{\mathcal F}$ is the Ideal
defining $\operatorname{supp}\mathcal F_{(r)}$, then $\mathcal F_{(r)}$
\src{258}
may be regarded as a Module on the subscheme $Y_{\mathcal F}$ of $X$ defined
by $\mathcal J_{\mathcal F}^{a}$. As in the preceding lemma, we may also
reduce to the case $\mathcal F=\mathcal F_{(r)}$, so that $\mathcal F$ comes
from a Module on $Y_{\mathcal F}$. The $Y_{\mathcal F}$ correspond to a
bounded family of quotient Modules of the $\mathcal O_{X_K}$, and hence arise
from a closed subscheme $Y$ of a scheme $X\times_S T$. We may then apply (2.5)
to $Y/T$ and $\mathcal L\otimes_XY$, which gives the conclusion.

We can now prove (2.2) by induction on the upper bound $r$ of the
$\dim\operatorname{supp}\mathcal F$. The statement is trivial for $r<0$;
suppose then that $r\geq0$ and that the statement has been proved for $r'<r$.
By (2.6), the $\mathcal F_{(r)}$ form a bounded set; hence, by (1.2) (i), so do
the kernels of the homomorphisms
$\mathcal L_K\longrightarrow\mathcal F_{(r)}$. There is therefore a coherent
Module $\mathcal L'$ on $X$ such that the kernels in question, and consequently
also the
\[
N_r(\mathcal F)=\operatorname{Ker}(\mathcal F\longrightarrow\mathcal F_{(r)}),
\]
are quotients of the Modules $\mathcal L'_K$. Since the $\mathcal F_{(r)}$ are
bounded, the $P_{\mathcal F_{(r)}}$ remain bounded, and the formula
\[
P_{\mathcal F}=P_{\mathcal F_{(r)}}+P_{N_r}
\]
shows that the $P_{N_r}$ satisfy the same condition $(b_s)$ as the
$P_{\mathcal F}$. Thus the $N_r$ satisfy conditions (a) and $(b_s)$, and by
the induction hypothesis the $(N_r)_{(s)}$ remain bounded. Now
$\mathcal F_{(s)}$ is an extension of $\mathcal F_{(r)}$ by
$(N_r)_{(s)}$; hence by (1.2) (ii), the $\mathcal F_{(s)}$ remain bounded. For the final assertion of (2.2), carry boundedness and the next-coefficient lower bound simultaneously in the induction. The bounded kernels of the fixed-source maps to $\mathcal F_{(s)}$ provide a single coherent source surjecting onto each $N_s$, but (1.2) (i) alone does not make the quotients $N_s$ bounded. Reapply the coefficient assertion inductively to those quotients; Hilbert-polynomial additivity and Lemma 2.6 then give the required lower bound for the next coefficient.
This completes the proof.

\fsection{3}{Hilbert schemes: definition and existence theorem}

Let $X$ be a prescheme over another prescheme $S$, and let $\mathcal F$ be a
quasi-coherent Module on $X$. Denote by
\[
\operatorname{Quot}(\mathcal F/X/S)
\]
the set of quasi-coherent quotient Modules of $\mathcal F$ that are flat over
$S$. Let now $S'\longrightarrow S$ be a base-change morphism, put
$X'=X\times_S S'$ and
$\mathcal F'=\mathcal F\otimes_{\mathcal O_S}\mathcal O_{S'}$. Thus $X'$ is
a prescheme over $S'$ equipped with a Module
\src{259}
$\mathcal F'$, and we may consider
$\operatorname{Quot}(\mathcal F'/X'/S')$. Put
\[
\underline{\operatorname{Quot}}_{\mathcal F/X/S}(S')
=\operatorname{Quot}(\mathcal F'/X'/S')
\qquad\text{(where $X'=X\times_SS'$)}.
\]
If now $S''\longrightarrow S'$ is an $S$-morphism, then
$X''=X\times_SS''$ is isomorphic to $X'\times_{S'}S''$, and $\mathcal F''$ to
$\mathcal F'\otimes_{\mathcal O_{S'}}\mathcal O_{S''}$. Since the inverse-image
functor
\[
\mathcal G'\longmapsto
\mathcal G'\otimes_{\mathcal O_{S'}}\mathcal O_{S''}
\]
from the category of quasi-coherent Modules on $X'$ to the category of
quasi-coherent Modules on $X''$ is right exact and transforms\footnote{The printed French omits the article \emph{les} before its direct object; the English was already grammatical.} $S'$-flat
Modules into $S''$-flat Modules, there is a natural map
\[
\operatorname{Quot}(\mathcal F'/X'/S')\longrightarrow
\operatorname{Quot}(\mathcal F''/X''/S''),
\]
so $\underline{\operatorname{Quot}}_{\mathcal F/X/S}(S')$ is a contravariant
functor of $S'$ (a prescheme over $S$), with values in the category of sets.
From now on we shall suppose that $X$ is projective over a noetherian $S$ and
that $\mathcal F$ is coherent; for simplicity we shall restrict to $S'$ over
$S$ that are locally noetherian.

\result{THEOREM}{3.1} Under these hypotheses, the contravariant functor
$\underline{\operatorname{Quot}}_{\mathcal F/X/S}$ on the category of locally
noetherian $S$-preschemes is representable by an $S$-prescheme
$\operatorname{Quot}_{\mathcal F/X/S}$ that is a disjoint union of a sequence
of projective $S$-schemes (in particular,
$\operatorname{Quot}_{\mathcal F/X/S}$ is locally of finite type over $S$).

We shall obtain such a decomposition as follows. Let $\mathcal O_X(1)$ be an
invertible sheaf on $X$ that is very ample relative to $S$. For every
polynomial $P(n)$ with rational coefficients, let
$\operatorname{Quot}^P(\mathcal F/X/S)$ be the subset of
$\operatorname{Quot}(\mathcal F/X/S)$ consisting of the coherent quotients
$\mathcal G$ of $\mathcal F$ that are flat over $S$ and whose Hilbert
polynomial at every $s\in S$ equals $P$. Put
\[
\underline{\operatorname{Quot}}^P_{\mathcal F/X/S}(S')
=\operatorname{Quot}^P(\mathcal F'/X'/S').
\]
This gives a subfunctor of
$\underline{\operatorname{Quot}}_{\mathcal F/X/S}$. The invariance property of
Hilbert polynomials recalled in Section 2 implies the following: the functor
$\underline{\operatorname{Quot}}_{\mathcal F/X/S}$ is representable if and
only if the $\underline{\operatorname{Quot}}^P_{\mathcal F/X/S}$ are
representable; in that case the representing $S$-prescheme
$\operatorname{Quot}_{\mathcal F/X/S}$ is isomorphic to the disjoint union of
the $\operatorname{Quot}^P_{\mathcal F/X/S}$ representing the functors
$\underline{\operatorname{Quot}}^P_{\mathcal F/X/S}$. Thus (3.1) follows from
the following theorem:
\src{260}
\result{THEOREM}{3.2} With the preceding notation, the functor
$\underline{\operatorname{Quot}}^P_{\mathcal F/X/S}$ is represented by a
projective $S$-prescheme $\operatorname{Quot}^P_{\mathcal F/X/S}$.

The remainder of this section is devoted to the proof of (3.2).

Let $\nu$ be an integer. For every $S'$ over $S$, denote by $A_\nu(S')$ the
set of quotients $\mathcal G=\mathcal F'/\mathcal K$ of
$\mathcal F'=\mathcal F\otimes_{\mathcal O_S}\mathcal O_{S'}$ that are
coherent and flat over $S'$, and satisfy the following conditions:
\begin{align*}
\text{a.}\quad &R^if'_*(\mathcal G(n))=0 &&\text{for $i>0$ and $n\geq\nu$},\\
\text{b.}\quad &R^if'_*(\mathcal K(n))=0 &&\text{for $i>0$ and $n\geq\nu$},\\
\text{c.}\quad &f'_*(\mathcal K(\nu+k))
=\mathcal S'_k f'_*(\mathcal K(\nu)) &&\text{for $k\geq0$}.
\end{align*}
For the last relation, $X$ is assumed to be written as the homogeneous prime
spectrum of a quasi-coherent graded algebra $\mathcal S_*$ over $S$, positively
graded and generated by $\mathcal S_1$, and we put
$\mathcal S'_* =\mathcal S_*\otimes_{\mathcal O_S}\mathcal O_{S'}$, so that
$X'$ is the homogeneous prime spectrum of $\mathcal S'_*$. To prove (3.2), one
readily reduces to the case $X=\mathbf P^r_S$ (since $S$ is a union of open
sets $U$ such that $X|U$ is a closed subscheme of $\mathbf P^r_U$ and
$\mathcal O_X(1)$ is induced by $\mathcal O_{\mathbf P^r_U}(1)$), and to the
case in which $\mathcal F$ has the form
$\mathcal O_{\mathbf P^r_S}^{N}$ and is therefore flat over $S$. In the
preceding discussion the sheaves $\mathcal K$ are then also flat over $S'$.
The Künneth relations ([5], III, \S~7) show that conditions (a) and (b) are
stable under base change and imply that, for $n\geq\nu$, the formation of
$f'_*(\mathcal G(n))$ and $f'_*(\mathcal K(n))$ commutes with base extension
(loc. cit.); hence (c) is also stable under base extension.

In other words, $A_\nu(S')$ is a contravariant functor of $S'$, more precisely
a subfunctor of
$A(S')=\underline{\operatorname{Quot}}^P_{\mathcal F/X/S}(S')$. As $\nu$
varies, one obtains an increasing sequence of subsets $A_\nu(S')$ of $A(S')$
whose union is $A(S')$, by a well-known theorem of Serre [5], III, \S~2. Now
suppose that $\mathcal G$ is a coherent quotient of $\mathcal F'$, flat over
$S'$, and that $s$ is a point of $S'$ such that base change
$\operatorname{Spec}(k(s))\longrightarrow S'$ gives a quotient
$\mathcal G_s$ of $\mathcal F_s$ satisfying conditions (a), (b), and (c), that
is, belonging to $A_\nu(\operatorname{Spec}(k(s)))$. Then there exists an open
neighborhood $U$ of $s$ on which the same conditions hold for
$\mathcal G|f'^{-1}(U)$, that is, this quotient belongs to $A_\nu(U)$. For (a)
and (b), this follows from the “theorem on holomorphic functions” [5], III,
\S~7; condition (c) follows from Nakayama's lemma and from the fact that in any
case one knows
$f'_*(\mathcal K(n+k))=\mathcal S'_k f'_*(\mathcal K(n))$
\src{261}
for $n$ sufficiently large and $k\geq0$ (cf. [5], III, \S~2). From these
observations one concludes the following (compare [4], IV). The functor $A$ is
representable if and only if the functors $A_\nu$ are representable; in that
case the $S$-prescheme $Q$ representing $A$ is an increasing union of open
subschemes $Q_\nu$ representing the $A_\nu$.

Let
\[
M_* =\bigoplus_{n\geq0}f_*(\mathcal F(n))=\mathcal S_*^{N},
\]
so that
\[
M'_* =M_*\otimes_{\mathcal O_S}\mathcal O_{S'}
=\bigoplus_{n\geq0}f'_*(\mathcal F'(n))=\mathcal S'_*{}^{N}.
\]
It follows from (a) that

\noindent (a') $f'_*(\mathcal G(n))$ is locally free of rank $P(n)$ for
$n\geq\nu$, [5], III, \S~7; and from (b), with $i=1$, that

\noindent (a'') $f'_*(\mathcal G(n))$ is a quotient Module of $M'_n$.

Moreover, knowledge of this quotient Module for $n=\nu$ determines, by (c),
the sub-Modules $f'_*(\mathcal K(n))$ of $M'_n$ for $n\geq\nu$, and therefore
determines $\mathcal K$ and hence $\mathcal G$. Thus one obtains an injective
map
\[
A_\nu(S')\longrightarrow\operatorname{Grass}_{P(\nu)}(M'_\nu)
\]
from $A_\nu(S')$ to the set of locally free quotient Modules of $M'_\nu$ of
rank $P(\nu)$, and hence a functorial homomorphism
\[
i_\nu:A_\nu(S')\longrightarrow
\underline{\operatorname{Grass}}_{P(\nu)}(M_\nu)(S'),
\]
where the functor on the right is represented by the Grassmann scheme
$\operatorname{Grass}_{P(\nu)}(M_\nu)$ (compare [4], V), which is projective
over $S$. I claim:

\result{LEMMA}{3.3} The functor $A_\nu(S')$ is representable, and the morphism
$Q_\nu\longrightarrow\operatorname{Grass}_{P(\nu)}(M_\nu)$ representing the
homomorphism $i_\nu$ is an immersion (which implies that $Q_\nu$ is
quasi-projective over $S$).

This assertion is equivalent to the following (compare [4], IV). Suppose given
a quotient Module $N$ of $M'_\nu$, locally free of rank $P(\nu)$. Then there
exists a subprescheme $Z$ of $S'$ such that for every locally noetherian
prescheme $T'$
\src{262}
over $S'$, the inverse image of $N$ on $T'$ belongs to
$\operatorname{Im}A_\nu(T')$ if and only if $T'\longrightarrow S'$ factors
through the subprescheme $Z$. Changing notation, we may suppose $S'=S$; that
is, we are given a quotient $N_\nu$ of $M_\nu$ by a sub-Module $R_\nu$. It
comes from an element of $A_\nu(S)$ if and only if it satisfies the following
two conditions:

\medskip
\noindent (i) $M_{\nu+k}/\mathcal S_kR_\nu$ is locally free of rank
$P(\nu+k)$ for $k\geq0$.

\smallskip
\noindent (ii) The subsheaf $\mathcal K$ of $\mathcal F$ defined by the graded
sub-Module
\[
R_* =\sum_{k\geq0}\mathcal S_kR_\nu
\]
of $M_*$ ([5], II, \S~3), and the quotient
$\mathcal G=\mathcal F/\mathcal K$, satisfy conditions (a) and (b) above.

These conditions are manifestly necessary. Conversely, if they hold, then the
sheaf $\mathcal G$ defined in (ii), being isomorphic to the sheaf associated
with the graded $\mathcal S_*$-Module $N_*$ that is the sum of the
$M_{\nu+k}/\mathcal S_kR_\nu$, is flat over $S$ (because its fibers are direct
factors of localizations of $N_*$\footnote{The printed text has $N$; homogeneous localization applies to the graded $\mathcal S_*$-Module $N_*$.} at homogeneous prime ideals of
$\mathcal S_*$), and has Hilbert polynomial $P$ by (i). In view of (ii), (a')
and (a'') hold, so for $n\geq\nu$ the natural homomorphism
$N_n\longrightarrow f_*(\mathcal G(n))$ is a surjective homomorphism between
locally free Modules of the same rank, and hence an isomorphism. Therefore
\[
f_*(\mathcal K(\nu+k))=\mathcal S_kR_\nu
\quad\text{for every $k\geq0$}.
\]
This proves that $\mathcal G\in A_\nu(S)$ and that $R_\nu$ is the element of
$\operatorname{Grass}_{P(\nu)}(M_\nu)$ defined by $\mathcal G$.

Criterion (i)--(ii) also applies after a base change
$S'\longrightarrow S$. We shall first prove that condition (i) holding after
the base change $S'\longrightarrow S$ can be expressed by requiring
$S'\longrightarrow S$ to factor through a certain subprescheme $Z$ of $S$.
Once this is done, replacing $S$ by $Z$ reduces us to the case in which (i)
already holds on $S$; since it is stable under base change, it remains to
express condition (ii). Let $U$ be the set of $s\in S$ such that the
cohomology in dimension $>0$ of the sheaves induced on the fiber $X_s$ by
$\mathcal G(n)$ and $\mathcal K(n)$ vanishes for $n\geq\nu$. We have already
observed that $U$ is open, and condition (ii) holds after a base change
$S'\longrightarrow S$ if and only if $S'\longrightarrow S$ factors through
$U$. This proves Lemma 3.3. It remains to prove the following lemma:

\result{LEMMA}{3.4} Let $S$ be a locally noetherian prescheme equipped with a
positively graded quasi-coherent algebra $\mathcal S_*$ generated by
$\mathcal S_1$; let $M_*$ be a quasi-coherent graded $\mathcal S_*$-Module of
finite type, $P$ a polynomial with rational coefficients, and $\nu$ an
integer. Then there exists an obviously unique subprescheme $Z$ of $S$
\src{263}
with the following property: for every prescheme $S'$ over $S$, the Module
$M_n\otimes_{\mathcal O_S}\mathcal O_{S'}$ is locally free of rank $P(n)$ for
every $n\geq\nu$ if and only if $S'\longrightarrow S$ factors through $Z$.

We may clearly suppose $S$ affine, hence noetherian. Then:

\result{LEMMA}{3.5} For every integer $N\geq\nu$, let $U_N$ be the open subset
of $S$ consisting of those $s\in S$ such that
\[
\operatorname{rank}_{k(s)}
M_{ns}\otimes_{\mathcal O_{S,s}}k(s)\leq P(n)
\quad\text{for every $\nu\leq n\leq N$}.
\]
Then the decreasing sequence of open subsets $U_N$ is stationary.

It is known [3], IV, that $S$ has a finite partition into reduced subschemes
$S_i$ such that each $M\otimes_{\mathcal O_S}\mathcal O_{S_i}$ is flat over
$S_i$. We may therefore suppose $M$ flat, and hence the $M_n$ flat. Finally,
we may clearly suppose $S$ connected. There then exist an integer $n_0$ and a
polynomial $Q$ such that
\[
\operatorname{rank}_{k(s)}
M_{ns}\otimes_{\mathcal O_{S,s}}k(s)=Q(n)
\quad\text{for $n\geq n_0$},
\]
[5], III, \S~7. First suppose $P<Q$\footnote{The printed text has
$P\neq Q$; the 1962 erratum prescribes $P<Q$, replaces the “$P=Q$” case by
the complementary case $P(n)\geq Q(n)$ for large $n$, and replaces the strict
inequalities in (*) by non-strict ones.}, so that $P(n)<Q(n)$ for large $n$.
Then plainly $U_N=\varnothing$ for $N$ large, and the sequence $U_N$ is
stationary. In the complementary case, $P(n)\geq Q(n)$ for large $n$, so
$U_N=U_{n_0}$ for $N\geq n_0$, and again the sequence $U_N$ is stationary.

In particular, the set $U_\infty$ of $s\in S$ such that
\[
\tag{*}
\operatorname{rank}_{k(s)}
M_{ns}\otimes_{\mathcal O_{S,s}}k(s)\leq P(n)
\quad\text{for every $n\geq\nu$}
\]
is open, since it is the intersection of the $U_N$. For the proof of (3.4), we
may therefore replace $S$ by the open set $U_\infty$\footnote{The printed text has $U$; the open subset just defined is $U_\infty$.}, reducing to the case in which
inequality (*) holds at every $s\in S$. Moreover:

\result{LEMMA}{3.6} Let $M$ be a Module on a locally noetherian prescheme $S$,
and let $r$ be an integer. Then there exists an obviously unique subprescheme
$Z$ of $S$ with the following property: for every $S'$ over $S$, the Module
$M\otimes_{\mathcal O_S}\mathcal O_{S'}$ is locally free of rank $r$ if and
only if $S'\longrightarrow S$ factors through $Z$. If
\[
\operatorname{rank}_{k(s)}
M_s\otimes_{\mathcal O_{S,s}}k(s)\leq r
\quad\text{for every $s$},
\]
then $Z$ is a closed subprescheme of $S$ (with $M$ assumed coherent).
\src{264}
Indeed, the preceding argument reduces us to the case in which inequality (*)
holds for every $s\in S$ (replacing $S$, if necessary, by the open subset of
points where it holds). We may then suppose that $M$ occurs in an exact
sequence
\[
\mathcal O_S^q\longrightarrow\mathcal O_S^r\longrightarrow M\longrightarrow0.
\]
The condition under consideration for $S'$ over $S$ also means that in the
corresponding exact sequence
$\mathcal O_{S'}^q\longrightarrow\mathcal O_{S'}^r
\longrightarrow M'\longrightarrow0$, the second arrow is an isomorphism; that
is, the first is zero. Hence the closed subprescheme $Z$ of $S$ defined by the
Ideal generated by the entries of the matrix defining
$\mathcal O_S^q\longrightarrow\mathcal O_S^r$ has the required property.

Return now to the proof of (3.4). Let $Z_n$ be the closed subscheme of $S$
associated by (3.6) with the Module $M_n$ and the integer $r=P(n)$, and let
$Z'_N$ be the infimum of the $Z_n$ for $\nu\leq n\leq N$. The $Z'_N$ form a
decreasing sequence of closed subschemes of $S$\footnote{The printed text has $Z$, which is defined only in the following sentence; each $Z'_N$ is constructed as a closed subscheme of $S$.}, and hence are necessarily
stationary. Let $Z$ be their constant value for $N$ large. This is the $Z$
sought in (3.4). This completes the proof of (3.4), and hence also of (3.3).

We have therefore proved that
$\underline{\operatorname{Quot}}^P_{\mathcal F/X/S}$ is represented by an
$S$-prescheme $Q$ that is an increasing union of open subschemes $Q_\nu$
quasi-projective over $S$. To go further, invoke Theorem 2.1, from which it
follows readily that $Q$ is quasi-compact (for it is the image of the
prescheme $S'$ of finite type over $S$ parametrizing the family of quotient
sheaves of the $\mathcal F_K$ with Hilbert polynomial $P$). Thus $Q$ equals one
of the $Q_\nu$, and so is quasi-projective over $S$. To prove that it is
projective over $S$, it remains to prove that it is proper over $S$. For this,
it suffices to invoke the valuative criterion for properness in the form [5],
II, 7.3.8. We need only verify:

\result{LEMMA}{3.7} Let $S$ be the spectrum of a discrete valuation ring, let
$s$ be its generic point, let $X$ be a prescheme over $S$, let $\mathcal F$ be
a quasi-coherent Module on $X$, and let $\mathcal G_s$ be a quasi-coherent
quotient Module of
$\mathcal F_s=\mathcal F\otimes_{\mathcal O_S}k(s)$ on $X_s$. Then there
exists a unique quasi-coherent quotient Module $\mathcal G$ of $\mathcal F$,
flat over $S$, whose restriction to $X_s$ is $\mathcal G_s$.
\src{265}
Indeed, if $\mathcal G_s=\mathcal F_s/\mathcal K_s$, it is enough to take the
largest subsheaf $\mathcal K$ of $\mathcal F$ inducing $\mathcal K_s$, [5], I,
9.4.2, and put $\mathcal G=\mathcal F/\mathcal K$. One checks readily that
this sheaf has the required properties.

Theorem 3.2, and consequently (3.1), is now completely proved.

The proof establishes at the same time:

\result{PROPOSITION}{3.8} Under the hypotheses of (3.2), let
$Q=\operatorname{Quot}^P_{\mathcal F/X/S}$, $X_Q=X\times_SQ$, and
$\mathcal F_Q=\mathcal F\otimes_{\mathcal O_S}\mathcal O_Q$; and let
$\mathcal G$ be the coherent quotient of $\mathcal F_Q$, flat over $Q$, with
relative Hilbert polynomial $P$, such that $(Q,\mathcal G)$ represents the
functor $\underline{\operatorname{Quot}}^P_{\mathcal F/X/S}$. There exists an
integer $\nu$ such that for $n\geq\nu$, $(f_Q)_*(\mathcal G(n))$ is a Module on
$Q$, locally free of rank $P(n)$ and very ample relative to $S$; that is, it
defines an immersion of $Q$ into a Grassmann scheme
$\operatorname{Grass}_{P(n)}(M)$ over $S$. In particular, for $n\geq\nu$, the
sheaf
\[
\bigwedge^{P(n)}(f_Q)_*(\mathcal G(n))
\]
on $Q$ is invertible and very ample relative to $S$.

As in (3.2), we may reduce to the case in which $\mathcal F$ is flat over $S$;
then it is enough to choose an integer $\nu$ such that $A_\nu=A$ (with the
preceding notation).

The most important application of (3.2) is the case
$\mathcal F=\mathcal O_X$. We then write
\[
\operatorname{Quot}_{\mathcal O_X/X/S}=\operatorname{Hilb}_{X/S},
\qquad
\operatorname{Quot}^P_{\mathcal O_X/X/S}=\operatorname{Hilb}^P_{X/S},
\]
and obtain a decomposition
\[
\operatorname{Hilb}_{X/S}=\coprod_P\operatorname{Hilb}^P_{X/S}.
\]
By definition, $\operatorname{Hilb}_{X/S}$ represents the functor
$\underline{\operatorname{Hilb}}_{X/S}(S')$, the set of closed subschemes of
$X'=X\times_SS'$ that are flat over $S'$; and
$\operatorname{Hilb}^P_{X/S}$ represents the corresponding subfunctor of
closed subschemes with a given Hilbert polynomial $P$. These preschemes are
also called the \emph{Hilbert prescheme} of $X$ over $S$ and the
\emph{Hilbert prescheme of index $P$}, respectively. The terminology is
justified by the role played in the theory by Hilbert polynomials. Their
difference in nature from the classical Chow varieties (which parametrize
cycles rather than varieties) is of the same order as that between
\src{266}
the Chow ring of cycle classes on a variety and the ring of sheaf classes on
the variety (as it appears\footnote{The printed French has feminine « telle », although its antecedent \emph{anneau} is masculine; the English was already idiomatic.} in the Riemann--Roch theorem [1]). Indeed, when
$X=\mathbf P^r_S$ and $S$ is the spectrum of a field, knowledge of the Hilbert
polynomial of a coherent Module $\mathcal F$ on $X$ is also equivalent to
knowledge of the Chern classes of $\mathcal F$, or of the class of $\mathcal F$
in the ring of classes of coherent sheaves on $X$.

\result{REMARKS}{3.9} Note also that the construction of
$\operatorname{Quot}_{\mathcal F/X/S}$ and
$\operatorname{Quot}^P_{\mathcal F/X/S}$ was reduced to the case
$X=\mathbf P^r_S$, $\mathcal F=\mathcal O_X^N$, with $\mathcal O_X(1)$ the
usual very ample sheaf; more precisely, the general
$\operatorname{Quot}_{\mathcal F/X/S}$ occur as closed subschemes of these.
Since the formation of the $\operatorname{Quot}^P_{\mathcal F/X/S}$ is plainly
compatible with base changes $S'\longrightarrow S$, we are reduced further to
$S=\operatorname{Spec}(\mathbf Z)$. Ultimately, then, one is reduced to
studying the following projective schemes over $\mathbf Z$:
\[
\mathfrak Q^P_{r,N}=
\operatorname{Quot}^{P}_{(\mathcal O_{\mathbf P^r_{\mathbf Z}})^N/
\mathbf P^r_{\mathbf Z}/\operatorname{Spec}(\mathbf Z)},
\]
and, more particularly, the absolute Hilbert schemes
\[
\operatorname{Hilb}^P_r=\mathfrak Q^P_{r,1}.
\]
A more detailed study of these schemes, beginning with the determination of
their connected components (are they connected?),\footnote{The 1962 erratum
notes that Hartshorne proved the connectedness of the
$\operatorname{Hilb}^P_r$ over an algebraically closed field and determined
the pairs $(r,P)$ for which $\operatorname{Hilb}^P_r\neq\varnothing$.} and
their irreducible components (by Serre [7], there may be irreducible components
lying entirely over a nonzero prime $p$), would be very interesting. Recall
Weil's question whether the irreducible components of the fibers of
$\operatorname{Hilb}^P_r$ over $s\in\operatorname{Spec}(\mathbf Z)$
correspond to “regular” extensions of the prime field, that is, whether they
are “relatively connected.” These questions may be more approachable for
Hilbert schemes than for “Chow varieties.”

\fsection{4}{Variants}

\noindent a. Under the hypotheses of (3.1), let $U$ be an open subset of $X$,
and let $A'$ be the subfunctor of
$A=\underline{\operatorname{Quot}}_{\mathcal F/X/S}$ such that $A'(S')$ is the
set of quotient Modules $\mathcal G$ of $\mathcal F'$, flat over $S'$, whose
support is contained in $U'$. Plainly $A'$ is represented by an open subset of
the prescheme $\operatorname{Quot}_{\mathcal F/X/S}$ representing $A$. It
follows that Theorems (3.1) and (3.2)
\src{267}
remain valid when $X$ is assumed quasi-projective over $S$ instead of
projective over $S$, provided “projective” is also replaced by
“quasi-projective” in the conclusions and
$\underline{\operatorname{Quot}}_{\mathcal F/X/S}(S')$ now denotes the set of
coherent quotients $\mathcal G$ of $\mathcal F'$, flat over $S'$, whose support
is proper over $S'$.

\medskip
\noindent b. More generally, one may impose on the quotients $\mathcal G$ of
$\mathcal F'=\mathcal F\otimes_{\mathcal O_S}\mathcal O_{S'}$, flat over $S'$,
all kinds of additional natural conditions stable under base change, thereby
obtaining subfunctors of
$\underline{\operatorname{Quot}}_{\mathcal F/X/S}$ whose representability may
be sought. In many cases, the usual criterion proves that these functors are
again represented by open subsets of
$\operatorname{Quot}_{\mathcal F/X/S}$. This is so, in particular, when one
imposes one of the following additional properties:

\smallskip
\noindent $1^\circ$ The dimensions of the prime cycles associated with the
Modules $\mathcal G_{s'}$ ($s'\in S'$)\footnote{The printed text has $\mathcal G_s$ although the quantified point is $s'$.} induced on the fibers $X'_{s'}$ belong
to a given set of integers.

\smallskip
\noindent $2^\circ$ (When $\mathcal F=\mathcal O_X$, so that $\mathcal G$
corresponds to a closed subprescheme $Y$ of $X'$): $Y$ is a smooth prescheme
([3], IV) over $S'$\footnote{In the six occurrences of this condition after base change, the printed text incorrectly retains $S$ or $X$; the family lives on $X'=X\times_SS'$ relative to $S'$.}, respectively normal over $S'$ (that is, the fibers $Y_{s'}$
are normal “over $k(s')$,”\footnote{The printed text has $k(s)$; the residue field of the indexed point $s'$ is $k(s')$.} meaning normal after every extension of the base
field), respectively (when $X$ is flat over $S$) an $S'$-local complete
intersection\footnote{The printed text has “$k$-local complete intersection”; the intermediate correction is relative to the base, and requires the base-changed $S'$ here.} in $X'$ over $S'$ (that is, the fibers $Y_{s'}$ are local complete
intersections in the $X'_{s'}$).

Other conditions would involve cohomological properties of the Modules
$\mathcal G_{s'}$ induced on the $X'_{s'}$, etc. Of course, a conjunction of
conditions each represented by an open subset $U_i$ of
$\operatorname{Quot}_{\mathcal F/X/S}$ is represented by the intersection of
these open subsets. For example, assigning to every $S'$ over $S$ the set of
closed subschemes $Y$ of $X'=X\times_SS'$ that are étale coverings ([3], I) of
$S'$ of a given rank $r$ gives a representable contravariant functor of $S'$.

\medskip
\noindent c. The preschemes $\operatorname{Hom}_S(X,Y)$,
$\Pi_{X/S}(Z/X)$, and $\operatorname{Isom}_S(X,Y)$ defined in [2], II, C,
no. 2 exist under suitable projectivity hypotheses and occur as open subsets
of suitable Hilbert preschemes. Since
\[
\operatorname{Hom}_S(X,Y)=\Pi_{X/S}((X\times_SY)/X),
\]
the case of $\operatorname{Hom}_S(X,Y)$\footnote{The printed text has $\operatorname{Hom}_S(X,X)$ here, contrary to the preceding displayed identity.} reduces to that of
$\Pi_{X/S}(Z/X)$. For every $S'$ over $S$, the set of sections of
$Z'=Z\times_SS'$ over $X'=X\times_SS'$ is in bijective correspondence with
the set of subpreschemes $\Gamma$ of $Z'$\footnote{The printed text has $Z$; these graphs are graphs of sections of $Z'=Z\times_SS'$.} (necessarily closed if $Z$ is
separated over $X$) such that the morphism $\Gamma\longrightarrow X'$ induced
by $Z'\longrightarrow X'$ is an isomorphism. In this way,
\src{268}
when $X$ is flat and proper over $S$ and $Z$ is quasi-projective over $S$,
$\Pi_{X/S}(Z/X)$ exists and occurs as an open subprescheme of
$\operatorname{Hilb}_{Z/S}$. Hence when $X$ is projective and flat over $S$
and $Y$ is quasi-projective over $S$, $\operatorname{Hom}_S(X,Y)$ exists and
occurs as an open subprescheme of
$\operatorname{Hilb}_{(X\times_SY)/S}$. When $X$ and $Y$ are both projective
over $S$, it follows at once that $\operatorname{Isom}_S(X,Y)$ also exists and
is represented by an open subset of $\operatorname{Hom}_S(X,Y)$. Similarly,
when $X$ is flat and projective over $S$ and $Y$ is quasi-projective over $S$,
the $S$-prescheme $\operatorname{Imm}_S(X/Y)$ corresponding to the subfunctor
of $\operatorname{Hom}_S(X,Y)$ whose elements are the $S'$-homomorphisms
$X'\longrightarrow Y'$ that are immersions is also represented by an open
subset of $\operatorname{Hom}_S(X,Y)$.

Let $\mathcal L$\footnote{The printed French uses plural « Soient » before a singular subject with a parenthetical counterpart; the English was already singular.} (respectively, $\mathcal M$) be an invertible sheaf on $X$
(respectively, $Y$), very ample relative to $S$. Then
$\mathcal L\otimes_{\mathcal O_S}\mathcal M$ is a very ample sheaf on
$X\times_SY$ relative to $S$. Consequently, for every polynomial $P$ with
rational coefficients, $\operatorname{Hilb}^P_{(X\times_SY)/S}$ is defined
and is a quasi-projective prescheme over $S$. It therefore induces on
$\operatorname{Hom}_S(X,Y)$ an open-and-closed subset quasi-projective over
$S$, which we denote $\operatorname{Hom}_S(X,Y)^P$. Thus the sections of
$\operatorname{Hom}_S(X,Y)^P$ over $S$ are the $S$-morphisms
$g:X\longrightarrow Y$ such that for every integer $n$,
\[
\chi\!\left(\left(\mathcal L\otimes_{\mathcal O_X}g^*(\mathcal M)\right)^{
\otimes n}\right)=P(n).
\]

In this way one obtains generalizations of Matsusaka's theorem asserting that
the automorphisms of a “polarized” projective variety form an algebraic group.
Here the assertion has a plainly more precise meaning, since this group is
defined as the solution of a universal problem. Note also that over an
algebraically closed field, the automorphism group considered formerly is
obtained from the “true” one defined here by quotienting out the nilpotent
elements. This explains why the older constructions have little chance of
working over a nonperfect base field: the Ideal of nilpotent elements appearing
after an extension of the base field need not be “defined over $k$.” The same
observation applies to most of the old-style constructions as well.

\fsection{5}{Infinitesimal study of Hilbert schemes}

It follows from the next result:

\src{269}
\result{PROPOSITION}{5.1} Let $S$ be a prescheme, let $S_0$ be a subprescheme
defined by a quasi-coherent Ideal $\mathcal J$ of square zero, let $X$ be an
$S$-prescheme, and let $\mathcal F$ be a quasi-coherent Module on $X$. Put
$X_0=X\times_SS_0$ and
$\mathcal F_0=\mathcal F\otimes_{\mathcal O_S}\mathcal O_{S_0}$, and let
$\mathcal G_0=\mathcal F_0/\mathcal K_0$ be a quasi-coherent quotient Module of
$\mathcal F_0$, flat over $S_0$. For every open subset $U$ of $X$, let
$\mathcal E(U)$ be the set of quasi-coherent quotient Modules $\mathcal G$ of
$\mathcal F|U$, flat over $S$, such that
$\mathcal G\otimes_{\mathcal O_S}\mathcal O_{S_0}=\mathcal G_0$. As $U$
varies, the $\mathcal E(U)$ are thus the sections of a sheaf $\mathcal E$ on
$X$. The sheaf of groups
\[
\mathcal A=\operatorname{Hom}_{\mathcal O_{X_0}}
(\mathcal K_0,\mathcal G_0\otimes_{\mathcal O_{S_0}}\mathcal J)
\]
acts naturally on $\mathcal E$, which thereby becomes a sheaf “formally
principal homogeneous under $\mathcal A$” (that is, for every open $U$ in $X$,
$\mathcal E(U)$ is empty or is a principal homogeneous set under
$\mathcal A(U)$).

It follows that:

\result{COROLLARY}{5.2} Suppose that locally on $X$ there exists an extension
$\mathcal G$ of $\mathcal G_0$ to a quotient of $\mathcal F$ flat over $S$
(that is, the fibers of the sheaf $\mathcal E$ are nonempty). Then there is a
canonical obstruction class
\[
c(\mathcal G_0)\in H^1(X,\mathcal A)
\]
whose vanishing is necessary and sufficient for the existence of a global
extension $\mathcal G$ of $\mathcal G_0$ to a quotient of $\mathcal F$ flat
over $S$. If this class vanishes, the set $\mathcal E(X)$ of all possible
extensions is a principal homogeneous set under
\[
\mathcal A(X)=\operatorname{Hom}_{\mathcal O_{X_0}}
(\mathcal K_0,\mathcal G_0\otimes_{\mathcal O_{S_0}}\mathcal J).
\]
\footnote{In this last formula the printed text has $\mathcal O_X$ in place of
$\mathcal O_{X_0}$, although the immediately preceding definition of
$\mathcal A$ fixes $\mathcal O_{X_0}$.}
In particular, a global extension exists if $H^1(X,\mathcal A)=0$.

\result{COROLLARY}{5.3} Suppose that
$Q=\operatorname{Quot}_{\mathcal F/X/S}$ exists (cf. 4 (a)) -- for example,
that $X$ is quasi-projective over a locally noetherian $S$ and $\mathcal F$ is
coherent. Let $x\in Q$ correspond to a residue extension $K=k(x)$ of $k(s)$
($s\in S$):
\src{270}
thus $x$ is defined by a coherent quotient Module
$\mathcal G_0=\mathcal F_0/\mathcal K_0$ of the Module
$\mathcal F_0=\mathcal F_K$ on the $K$-prescheme $X_K$. Let $\mathcal A$ be the
coherent sheaf on $X_K$ defined by
\[
\mathcal A=\operatorname{Hom}_{\mathcal O_{X_0}}(\mathcal K_0,\mathcal G_0).
\]
Then the Zariski tangent space to the fiber $Q_s$ at $x$ (the $K$-dual of
$\mathfrak m/\mathfrak m^2$, where $\mathfrak m$ is the maximal ideal of
$\mathcal O_{Q_K,x}$) is canonically isomorphic to $H^0(X_K,\mathcal A)$.

The description of the Zariski tangent space admits the following
generalization. Given an $S$-morphism $g:S'\longrightarrow Q$, that is, a
section $g'$ of $Q'=Q\times_SS'$ over $S'$, the Module
\[
\Omega=g^*(\Omega^1_{Q/S})=g'^*(\mathcal J/\mathcal J^2)
\]
(where $\mathcal J$ is the Ideal on $Q'$ defined by the section $g'$ of $Q'$
over $S'$) is characterized by the following formula, functorial in the
coherent Module $\mathcal M$ on $S'$:
\[
\operatorname{Hom}_{\mathcal O_{S'}}(\Omega,\mathcal M)
\simeq H^0(X',\mathcal A),
\]
where $\mathcal A$ is the Module on $X'=X\times_SS'$ defined by
\[
\mathcal A=\operatorname{Hom}_{\mathcal O_{X'}}
(\mathcal K,\mathcal G\otimes_{\mathcal O_{S'}}\mathcal M).
\]
\footnote{The printed text puts $\mathcal O_S$ in this tensor product; here
$\mathcal G$ and $\mathcal M$ are Modules over $S'$, and the square-zero
deformation construction requires the tensor product over $\mathcal O_{S'}$.}
Here $\mathcal G=\mathcal F'/\mathcal K$ is the quotient Module of
$\mathcal F'=\mathcal F\otimes_{\mathcal O_S}\mathcal O_{S'}$ corresponding
to $g$. Indeed, apply (5.1) with $S_0$ replaced by $S'$ and $S$ by the
prescheme $D(\mathcal M)=(S',\mathcal O_{S'}+\mathcal M)$, where $\mathcal M$
is regarded as an Ideal of square zero.

When $\mathcal F=\mathcal O_X$ in (5.1), the datum of $\mathcal G_0$
corresponds to a closed subprescheme $Y_0$ of $X_0$, flat over $S_0$, defined
by the Ideal $\mathcal J_0=\mathcal K_0$. Formula (*) then gives
\[
\mathcal A=\operatorname{Hom}_{\mathcal O_{X_0}}
(\mathcal J_0/\mathcal J_0^2,
\mathcal O_{Y_0}\otimes_{\mathcal O_{S_0}}\mathcal J),
\]
where $\mathcal J_0/\mathcal J_0^2$ is interpreted as the conormal sheaf of $Y_0$
in $X_0$, also denoted
\src{271}
$\check{\mathcal N}_{Y_0/X_0}$. It is then useful to regard $\mathcal A$ as a
Module on $Y_0$ and compute $H^0$ and $H^1$ on $Y_0$. Moreover, if $Y_0$ is
locally a complete intersection in $X_0$, with $X$ flat over $S$, then local
extensions in (5.1) are guaranteed to exist; furthermore,
$\mathcal J_0/\mathcal J_0^2$ is locally free on $Y_0$, and one may write
\[
\mathcal A=\mathcal N_{Y_0/X_0}\otimes_{\mathcal O_{S_0}}\mathcal J,
\]
where the first factor on the right is the normal sheaf of $Y_0$ in $X_0$.
Using the fundamental criterion of smoothness ([3], III, 3.1), one obtains for
example:

\result{COROLLARY}{5.4} Under the hypotheses of (5.3), suppose that
$\mathcal F=\mathcal O_X$, that $X$ is flat over $S$, and that the closed
subprescheme $Y_0$ of $X_0$ corresponding to $\mathcal G_0$ is locally a
complete intersection. Then the Zariski tangent space to $Q_s$ at $x$ is
canonically isomorphic to $H^0(Y_0,\mathcal N_{Y_0/X_0})$. If
$H^1(Y_0,\mathcal N_{Y_0/X_0})=0$, then the Hilbert prescheme
$\operatorname{Hilb}_{X/S}$ is smooth over $S$ at $x$.
($\mathcal N_{Y_0/X_0}$ is the normal sheaf of $Y_0$ in $X_0$.)

\result{REMARK}{5.5} This statement applies in particular when $Y_0$ is a
complete intersection in $X_0$ defined by one equation, that is, a positive
“Cartier divisor.” Then $\mathcal N_{Y_0/X_0}$ is isomorphic to the sheaf on
$Y_0$ induced by the invertible sheaf $\mathcal J_0^{-1}$ on $X_0$ defined by the
divisor $Y_0$. This situation occurs especially in the study of families of
positive divisors on a nonsingular projective variety $X_0$. The isomorphism
between the Zariski tangent space at the point $x$ of $Q$ (or, equivalently,
of the open subset $D$ of $Q$ corresponding to divisors) and
$H^0(Y_0,\mathcal N_{Y_0/X_0})$ was known in classical algebraic geometry as
the “characteristic homomorphism” from the former to the latter. It was defined
only when $x$ was a smooth point of the parameter variety $T$ of a “complete
continuous family” of divisors, that is, from our point of view, of an
irreducible component of the scheme $D$ endowed with its induced reduced
structure. The
tangent space to $T$ at $x$ is then a subspace of the tangent space to $D$ at
$x$, so the classical characteristic homomorphism is indeed injective, but it
is surjective only under additional conditions, for example when $D$ is
integral at $x$. In fact, Zappa [8] constructed an example
\src{272}
(with $X$ a nonsingular projective surface over the complex numbers) in which
the characteristic homomorphism is not surjective even at the generic point of
$T$. Thus $D$ is not integral even at the generic point of the irreducible
component under consideration. This shows especially strikingly how varieties
with nilpotent elements are necessary for understanding phenomena of the most
classical theory of surfaces.\footnote{The 1962 erratum adds that Mumford had
just constructed an irreducible component of the Hilbert scheme of
$\mathbf P^3_{\mathbf C}$, whose general points represent nonsingular curves of
degree 14 and genus 24, that is nonreduced at its generic point; blowing up one
of these curves also gives a regular projective threefold whose formal moduli
scheme is generically nonreduced.}

\subsection*{5.6}

In (5.4) we gave a smoothness criterion applying in particular to schemes of
divisors. Kodaira gave a different criterion in [6], namely the vanishing of
$H^1(X_0,\mathcal L)$, where $\mathcal L=\mathcal J_0^{-1}$ is the invertible
sheaf on $X_0$ defined by the divisor $Y_0$. This criterion is valid when $S$
is the spectrum of a field of characteristic zero and was proved in [6] by
transcendental methods when the base field is $\mathbf C$. More generally, with
$S$ arbitrary again, Kodaira's condition is sufficient for the canonical
morphism
\[
D\longrightarrow\operatorname{Pic}_{X/S}
\]
from the prescheme of divisors to the Picard prescheme of $X/S$ to be smooth at
the point $x$ in question (as follows readily from the usual smoothness
criterion once the existence of $\operatorname{Pic}_{X/S}$ is known). If, in
addition, $\operatorname{Pic}_{X/S}$ is smooth over $S$ at the image of $x$
(for example, if $\operatorname{Pic}_{X/S}$ is smooth over $S$), then $D$ is
smooth over $S$ at $x$. On the other hand, Cartier proved that every group
prescheme locally of finite type over a field $k$ of characteristic zero is
smooth over $k$. Combining these two results recovers Kodaira's result. These
remarks also show that over a field $K$ of characteristic $p>0$, when
$\operatorname{Pic}_{X/S}$ is not smooth over $K$\footnote{The printed text has $k$ here; the base field of this positive-characteristic sentence is $K$.} (as when $X$ is the Igusa
surface), the condition $H^1(X_0,\mathcal L)=0$ implies, on the contrary, that
$D$ is not smooth at $x$, and is even nonreduced at $x$ if $K$ is algebraically
closed.

Finally, let us record the following result, which plays an important role in
the infinitesimal study of fiber spaces:

\result{PROPOSITION}{5.7} Let $X$ be a finite flat prescheme over a locally
noetherian $S$, and let $Z$ be a prescheme over $X$ such that
$\Pi_{X/S}(Z/X)$ exists (as it does if $Z$ is quasi-projective over $X$). If
$Z$ is smooth over $X$, then $\Pi_{X/S}(Z/X)$ is smooth over $S$.

This follows immediately from the definition and the usual smoothness
criterion ([3], III, 3.1). Note that when $X$ is finite and flat over $S$, the
existence of $\Pi_{X/S}(Z/X)$\footnote{The printed text has $Z/K$ here; no $K$
is defined, and the section functor under discussion since (5.7) is
$\Pi_{X/S}(Z/X)$.} can be treated very elementarily, without
\src{273}
using the theory of Hilbert schemes. For example, if $X$ is radicial over $S$,
then $\Pi_{X/S}(Z/X)$ exists without any restriction on $Z$. As another
example, let $T$ be an $S$-prescheme and let $T_n$ be the “infinitesimal
neighborhood of order $n$” of the diagonal of $T\times_ST$ in $T\times_ST$,
equipped with the morphisms $p_1,p_2:T_n\longrightarrow T$ induced by the two
projections. Regard $T_n$ as a prescheme finite over $T$ via $p_1$, and suppose
in addition that it is flat over $T$ (as it is when $T$ is smooth over $S$).
For every prescheme $X$ over $T$, put
\[
(X/T/S)^{(n)}=
\Pi_{T_n/T}\bigl(p_2^*(X/T)/T_n\bigr).
\]
\footnote{The printed text places $T_n/S$ under the symbol $\Pi$. Proposition
5.7 is applied here to the finite flat morphism $p_1:T_n\to T$, and the result
is explicitly a prescheme over $T$; hence the denominator is $T_n/T$.}
This is a prescheme over $T$ called the \emph{bundle of germs of sections of
order $n$ of $X$ over $T$} (relative to $S$). It depends functorially on $X$
and is smooth over $T$ if $X$ is.

\fsection{6}{Relations with the norm and symmetric products}

Let $S$ be a prescheme, let $X$ and $Y$ be $S$-preschemes, and let
\[
u:(X/S)^n\longrightarrow Y
\]
be a symmetric $S$-morphism from the $n$-th Cartesian power of $X/S$ to $Y$.
For simplicity, suppose $S$ locally noetherian and $X$ and $Y$ of finite type
over $S$. To every coherent Module $\mathcal F$ on $X$ whose support is finite
over $S$, which is flat over $S$ and has relative rank $n$ over $S$ (that is,
$f_*(\mathcal F)$ is a locally free Module of rank $n$ on $S$), one can then
naturally associate a section of $Y$ over $S$:
\[
\pi^u_{X/S}(\mathcal F)\in\Gamma(Y/S).
\]
We shall not give the formal definition here, but merely note that the resulting
formalism is a natural generalization of the usual formalism of norms and
traces. For arbitrary separated $X/S$, the norm law has canonical target the divided-power space $\Gamma^n_S(X)$. If $X$ is flat over $S$ and the $n$-th symmetric power exists (for example, when the orbits of $\mathfrak S_n$ on $(X/S)^n$ lie in affine open subsets), the canonical map $\operatorname{Symm}^n_S(X)\to\Gamma^n_S(X)$ is an isomorphism, and one obtains a canonical element
\[
\pi_{X/S}(\mathcal F)\in\Gamma(\operatorname{Symm}^n_S(X)/S),
\]
\src{274}
from which the $\pi^u_{X/S}(\mathcal F)$ can be recovered. Another important
case is that in which $X$ is a commutative monoid over $S$, $X=Y$, and $u$
comes from the composition law of $X$. We then write simply $\pi(\mathcal F)$
for the section of $X$ over $S$ associated with the Module $\mathcal F$ on
$X$.

Suppose now that $\mathcal F$ is a coherent Module on $X$ such that
$\operatorname{Quot}_{\mathcal F/X/S}$ exists, or at least such that the
functor $\underline{\operatorname{Quot}}^n_{\mathcal F/X/S}$ assigning to every
$S'$ over $S$ the set of coherent quotient sheaves $\mathcal M$ of
$\mathcal F'=\mathcal F\otimes_{\mathcal O_S}\mathcal O_{S'}$ that are flat
over $S'$ and have relative rank $n$ is represented by an $S$-prescheme
$\operatorname{Quot}^n_{\mathcal F/X/S}$. (When $X$ is quasi-projective over
$S$, $\operatorname{Quot}^n_{\mathcal F/X/S}$ does exist and, in the notation
of Section 3, is $\operatorname{Quot}^P_{\mathcal F/X/S}$, where $P$ is the
constant polynomial $n$.) Since the formation of
$\pi^u_{X/S}(\mathcal M)$ is compatible with base change, one obtains a
canonical morphism
\[
\pi^u_{X/S}:\operatorname{Quot}^n_{\mathcal F/X/S}\longrightarrow Y,
\]
and, in general, \[ \pi_{X/S}:\operatorname{Quot}^n_{\mathcal F/X/S} \longrightarrow\Gamma^n_S(X). \] When $X$ is flat over $S$, identify this target with $\operatorname{Symm}^n_S(X)$.
Independently of the preceding identification, when $X$ is separated over $S$, the finite-flat-family construction for $\mathcal F=\mathcal O_X$ gives a morphism
\[
\pi_{X/S}:\operatorname{Hilb}^n_{X/S}
\longrightarrow\operatorname{Symm}^n_S(X).
\]
This is plainly an isomorphism for $n=0$ or $n=1$. But for
$n>1$\footnote{The printed text repeats $n\geq1$, contradicting the
isomorphism just asserted for $n=1$.}, even when
$S$ is the spectrum of a field $k$ and $X$ is smooth over $S$, it is not in
general an isomorphism, or even an injective morphism: a zero-dimensional
subscheme of $X$ (corresponding, for example, to an ideal $\mathfrak i$
primary for the maximal ideal in a local ring $\mathcal O_{X,x}$, for a closed
point $x$ of $X$) is not determined by the cycle it defines (in this case, by
the codimension over $k$ of $\mathfrak i$ in $\mathcal O_{X,x}$). One can say
only the following, with $S$ arbitrary again:

\medskip
\noindent a. If $X$ is smooth over $S$, the norm morphism induces an
isomorphism from the open subset of $\operatorname{Hilb}^n_{X/S}$ classifying
étale coverings of rank $n$ contained in $X$ (cf. no. 4, (b)) onto the open
subset of $\operatorname{Symm}^n_S(X)$ corresponding to $n$-cycles without
multiple components.
\src{275}

\medskip
\noindent b. If, in addition, $X$ has relative dimension 1 over $S$, then the
norm morphism is an isomorphism from $\operatorname{Hilb}^n_{X/S}$ onto
$\operatorname{Symm}^n_{X/S}$.

This is because a zero-dimensional subscheme of a nonsingular algebraic curve
is determined by its corresponding cycle. The same observation applies more
generally to positive Cartier divisors on a nonsingular algebraic scheme (and
it is possible that in this very special case the Chow variety gives the same
object as the Hilbert scheme).

\fsection{7}{Complements and questions}

As J.-P. Serre observed, a well-known example of Nagata shows that one can find
a scheme $S$, the spectrum of a field $k$; an $S$-scheme $S'$, the spectrum of
a quadratic extension $k'$ of $k$; and a proper smooth (but nonprojective)
$S'$-scheme $X$ of dimension 3 such that $\Pi_{S'/S}(X/S')$ does not exist. It
follows a fortiori that the Hilbert scheme $\operatorname{Hilb}^2_{X/S}$ does
not exist (nor even the $k$-scheme that would represent the étale coverings of
$S$ of rank 2 contained in $X$, and a fortiori not the symmetric square of
$X$; cf. the preceding section). This places serious limitations on the
possibility of nonprojective constructions in algebraic geometry. (It is,
however, plausible that no such limitations arise in analytic geometry, any
more than they do in formal geometry (cf. [2], II).) On the other hand, if $X$
is a proper scheme over $S=\operatorname{Spec}k$ and $Z$ is quasi-projective
over $X$, then $\Pi_{X/S}(Z/X)$ exists and is a scheme that is a disjoint union
of a sequence of schemes quasi-projective over $S$ (as in the projective case
(3.1)). Indeed, one reduces to the case in which $X$ itself is projective by
dominating $X$ by a projective $S$-scheme $X'$. We shall not give the details
of the proof, which also uses the factorization result for a finite morphism
recorded in [2], I, A, no. 2 (b). The method succeeds because, with $S$ the
spectrum of a field, the $X'$ occurring in Chow's lemma is automatically flat
over $S$.

I do not know whether the result remains true without a hypothesis on $S$,
assuming only that $X$ is proper and flat over $S$ and $Z$ quasi-projective
over $X$. An important case in applications is when $Z$ is a closed subscheme
of $X$; if $\Pi_{X/S}(Z/X)$ exists, it is necessarily a closed subscheme of
$S$. It may be
\src{276}
constructed directly and quite simply when $X$ is projective over $S$, without
using the theory of Hilbert schemes. More generally, the same method shows that
if $Z$ is affine over $X$, then $\Pi_{X/S}(Z/X)$ exists and is affine over $S$.
It also shows that if $X$ is proper and flat over $S$ (not necessarily
projective over $S$), then for every locally trivial vector bundle $Z$ over
$X$, $\Pi_{X/S}(Z/X)$ exists and is a vector bundle over $S$. It would be
desirable to revisit and unify these results.

\section*{BIBLIOGRAPHY}
\addcontentsline{toc}{section}{Bibliography}

\begin{description}
\item[{[1]}] BOREL (A.) and SERRE (J.-P.). -- Le théorème de Riemann--Roch,
\emph{Bull. Soc. math. France}, vol. 86, 1958, pp. 97--136.

\item[{[2]}] GROTHENDIECK (Alexander). -- Technique de construction et
théorèmes d'existence en géométrie algébrique, I, II, III, \emph{Séminaire
Bourbaki}, vol. 13, 1959/60, no. 190, 29 pp., no. 195, 22 pp.; and vol. 14,
1960/61, no. 212, 20 pp.

\item[{[3]}] GROTHENDIECK (Alexander). -- Séminaire de Géométrie algébrique,
I, II, III, IV. -- Paris, Institut des hautes Études scientifiques, 1960/61
(mimeographed).

\item[{[4]}] GROTHENDIECK (Alexander). -- Technique de construction en
géométrie analytique, IV, V, \emph{Séminaire Cartan}, vol. 13, 1960/61, nos. 11
and 12.

\item[{[5]}] GROTHENDIECK (A.) and DIEUDONNÉ (J.). -- \emph{Éléments de
géométrie algébrique}, I: Le langage des schémas. -- Paris, Presses
universitaires de France, 1960 (Institut des hautes Études scientifiques,
Publications mathématiques, 4); II, III, IV (forthcoming).

\item[{[6]}] KODAIRA (K.). -- Characteristic linear systems of complete
continuous systems, \emph{Amer. J. of Math.}, vol. 78, 1956, pp. 716--744.

\item[{[7]}] SERRE (Jean-Pierre). -- Exemples de variétés projectives en
caractéristique $p$ non relevables en caractéristique 0, \emph{Proc. Nat. Acad.
Sc. U.S.A.}, vol. 47, 1961, pp. 108--109.

\item[{[8]}] ZAPPA (G.). -- Sull'esistenza, sopra le superficie algebriche, di
sistemi continui completi infiniti, la cui curva è a serie caratteristica
incompleta, \emph{Pont. Acad. Sc. Acta}, vol. 9, 1945, pp. 91--93.
\end{description}

\section*{ADDENDUM}
\addcontentsline{toc}{section}{Addendum}
\begin{center}
{\small [added during proof correction]}
\end{center}

It now appears that the conjectures of [2], III, no. 8 are false, even for
nonsingular varieties over a field of characteristic zero, both as regards the
existence and the quasi-projectivity of the quotient, and even when $G$ acts
with closed graph.

\par\smallskip
\begingroup\footnotesize
\endgroup
